% ------------------------------------------------------------
\escenario{ESC-05}{Un cliente alterado intenta hacer trampa}

\begin{tabla}{|L{3.2cm}|Y|}
\fichacab
\bloque{Trazabilidad}
\parte{Característica}{QA-03 Seguridad: integridad.}
\parte{Stakeholders}{STK-15 Especialista en ciberseguridad, STK-10 Equipo de pruebas, STK-05 Arquitecto de software.}
\parte{Concerns}{CRN-23: \emph{``Doy por hecho que alguien va a modificar el cliente, así que voy a probar con uno alterado. Necesito comprobar que no puede hacer jugadas ilegales ni ver información oculta del rival.''} CRN-10 (quién manda sobre el estado), CRN-11 (protocolo propio y versiones).}
\parte{Restricciones y dependencias}{CON-03 (WebSockets prohibidos, protocolo propio), REQ-01 (multijugador en línea).}
\parte{Tensiones y preguntas}{TEN-01 Latencia contra integridad de la partida.}
\parte{Tipo y prioridad}{Exploratorio, en tiempo de ejecución (bajo ataque). Prioridad (A, A).}
\bloque{Escenario}
\parte{Fuente del estímulo}{Especialista en ciberseguridad que actúa como atacante con un cliente modificado.}
\parte{Estímulo}{Envía al servidor, dentro y fuera de su turno: movimientos ilegales, barreras superpuestas o que encierran al rival, más barreras de las que le quedan, mensajes truncados o mal formados, mensajes de una versión anterior del protocolo y solicitudes de datos de la cuenta del rival (correo, dirección IP).}
\parte{Ambiente}{Partida en línea en curso contra un jugador que usa el cliente original, con el servidor atendiendo otras 20 partidas.}
\parte{Artefacto}{Servidor de partidas, gestión de turnos y estado de partida, y capa de red sobre TCP (protocolo propio).}
\parte{Respuesta}{El servidor rechaza todo lo que no es válido sin cambiar la partida, el rival no ve ninguna jugada ilegal, el cliente alterado no obtiene datos que no le corresponden y el intento queda registrado con la cuenta que lo hizo. Las demás partidas siguen funcionando.}
\parte{Medida de la respuesta}{Con una batería de 200 mensajes inválidos definida por STK-15, de los cuales 50 se generan de forma aleatoria (fuzzing): el 100~\% se rechaza; el estado de la partida es idéntico antes y después de cada mensaje; 0 respuestas contienen datos privados del rival; el servidor no se detiene ni se reinicia; las otras 20 partidas siguen cumpliendo ESC-01; y el 100~\% de los intentos queda registrado con cuenta y hora.}
\end{tabla}

\textbf{Por qué es relevante.} Es la prueba de aceptación de la decisión más cara de cambiar del proyecto: quién manda sobre el estado de la partida (CRN-10). Además, al estar prohibidos los WebSockets, el equipo escribe su propio protocolo sobre TCP, con su propio delimitado de mensajes y su propio manejo de versiones (CRN-11). Todo ese código nuevo es superficie de ataque, y los mensajes mal formados son justo lo que lo pone a prueba. El escenario también mide la tensión con el desempeño: la validación en el servidor no puede degradar a las demás partidas por debajo de ESC-01.

\textbf{Cómo se verifica.} Se construye un cliente de prueba que habla el protocolo directamente, sin interfaz, y se ejecuta la batería completa. Antes y después de cada mensaje se compara una huella del estado de la partida en el servidor. Se revisan las respuestas buscando campos de la cuenta del rival y, en paralelo, se miden los tiempos de respuesta de las otras 20 partidas.
